\section{Authorization}
\label{sec:Authorization}

Authorization is enforced in Quarkus using a server-side Policy Enforcement
Point (PEP) and a project-specific in-process Policy Decision Point (PDP).

\subsection{Identity modes}

Identity resolution is pluggable. The mode is selected via
\code{ARCHIMESH\_IDENTITY\_MODE} (or \code{app.identity.mode}):

\begin{description}
\item[\code{bootstrap}] Local/dev mode. No external identity provider.
  \begin{itemize}
  \item REST: identity from \code{X-Archimesh-User} and
    \code{X-Archimesh-Roles} headers.
  \item WebSocket: identity from query parameters \code{user} and
    \code{roles}.
  \item Intended for local development and simple standalone deployments.
  \end{itemize}

\item[\code{proxy}] Trusted forwarded headers from a reverse proxy.
  \begin{itemize}
  \item REST: identity from trusted forwarded headers (default
    \code{X-Forwarded-User}, \code{X-Forwarded-Roles}).
  \item WebSocket: identity from trusted forwarded headers captured during
    the handshake.
  \item Header names are configurable via
    \code{app.identity.proxy.user-header} and
    \code{app.identity.proxy.roles-header}.
  \item Do not expose \code{proxy} mode directly to untrusted clients
    without a trusted proxy in front.
  \end{itemize}

\item[\code{oidc}] Direct principal/role resolution through Quarkus auth.
  \begin{itemize}
  \item REST: identity from the authenticated Quarkus
    \code{SecurityContext}.
  \item WebSocket: identity from the authenticated WebSocket upgrade
    principal and role checks captured during the handshake.
  \item Can use local JWT verification or an external OIDC provider.
  \item No reverse proxy required in this mode.
  \end{itemize}
\end{description}

\subsection{Role model}

Role names from the identity provider are normalized into three canonical
internal roles:

\begin{description}
\item[\code{admin}] Catalog-wide administration: register/delete models,
  manage ACLs, run compactions, export/import, view diagnostics.
\item[\code{model\_reader}] Read model snapshots and join read-only
  sessions.
\item[\code{model\_writer}] Submit operations, acquire locks, send
  presence. Implies \code{model\_reader}.
\end{description}

Role names are configurable:

\begin{lstlisting}[caption={Role configuration}]
# Default canonical role names
app.authz.admin-role=admin
app.authz.reader-role=model_reader
app.authz.writer-role=model_writer

# Optional provider-specific aliases (comma-separated)
app.authz.admin-role-aliases=realm-admin,archimesh-admin
app.authz.reader-role-aliases=realm-viewer,archimesh-reader
app.authz.writer-role-aliases=realm-editor,archimesh-writer
\end{lstlisting}

Aliases normalize provider-specific group names (e.g. Keycloak realm roles,
Auth0 permissions) into the canonical internal roles before the PDP evaluates
policy.

\subsection{Policy model}

The PDP decides on these actions:

\begin{itemize}
\item Admin catalog actions (list models, create, delete)
\item Model join
\item Snapshot read
\item Submit operations
\item Lock acquire/release
\item Presence broadcast
\item Tag creation
\item ACL management
\item Export/import
\item Compaction
\end{itemize}

Authorization is deny-by-default when \code{ARCHIMESH\_AUTHZ\_ENABLED=true}.

\subsubsection{Model-scoped ACLs}

Per-model ACLs are stored on the \code{:Model} node as
\code{adminUsers}, \code{writerUsers} and \code{readerUsers} arrays.
When an ACL is configured for a model, it overrides the generic reader/writer
role checks for that model:

\begin{itemize}
\item A user listed in \code{adminUsers} gets admin access to that model.
\item A user listed in \code{writerUsers} can submit ops and lock for that
  model.
\item A user listed in \code{readerUsers} can read snapshots for that
  model.
\item Catalog-wide admin actions still require the global \code{admin}
  role.
\end{itemize}

\subsection{JWT and OIDC setup}

\subsubsection{Local JWT (dev mode)}

For local development with bearer tokens, use \code{app.identity.mode=oidc}
with Quarkus JWT validation:

\begin{lstlisting}[caption={Local JWT config}]
app.authz.enabled=true
app.identity.mode=oidc
mp.jwt.verify.publickey.location=src/test/resources/jwt/publicKey.pem
mp.jwt.verify.issuer=https://archimesh.dev
\end{lstlisting}

Mint tokens with the dev script:

\begin{lstlisting}[language=bash,caption={Mint dev JWT}]
./scripts/mint-dev-jwt.sh --user alice --roles admin,model_writer,model_reader
\end{lstlisting}

Use the token in requests:

\begin{lstlisting}[language=bash,caption={Bearer token usage}]
TOKEN=$(./scripts/mint-dev-jwt.sh --user alice --roles admin)
curl -s -H "Authorization: Bearer $TOKEN" \
  http://localhost:8081/admin/models | jq .
\end{lstlisting}

\subsubsection{External OIDC}

For production with an external provider (Keycloak, Auth0):

\begin{lstlisting}[caption={External OIDC config}]
app.authz.enabled=true
app.identity.mode=oidc
quarkus.oidc.enabled=true
quarkus.oidc.auth-server-url=https://idp.example/realms/archimesh
quarkus.oidc.client-id=archimesh-server
quarkus.oidc.application-type=service
quarkus.oidc.roles.role-claim-path=groups
\end{lstlisting}

Concrete provider examples are in \filepath{server/examples/}.

\subsection{Authorization diagnostics}

The endpoint \code{GET /admin/auth/diagnostics} returns the resolved
identity mode, subject, normalized roles and authorization status for the
current request. The admin UI exposes this through the \textbf{Auth Check}
action. This is useful for debugging misconfigured role mappings or proxy
headers.

\subsection{Configuration summary}

\begingroup\small
\begin{longtable}{p{46mm}p{54mm}p{54mm}}
\toprule
Property & Env var & Default \\
\midrule
\code{app.authz.enabled} & \code{ARCHIMESH\_AUTHZ\_ENABLED} &
  \code{false} \\
\code{app.identity.mode} & \code{ARCHIMESH\_IDENTITY\_MODE} &
  \code{bootstrap} \\
\code{app.identity.proxy.user-header} &
  \code{ARCHIMESH\_PROXY\_USER\_HEADER} &
  \code{X-Forwarded-User} \\
\code{app.identity.proxy.roles-header} &
  \code{ARCHIMESH\_PROXY\_ROLES\_HEADER} &
  \code{X-Forwarded-Roles} \\
\code{app.authz.admin-role} & \code{ARCHIMESH\_AUTHZ\_ADMIN\_ROLE} &
  \code{admin} \\
\code{app.authz.reader-role} & \code{ARCHIMESH\_AUTHZ\_READER\_ROLE} &
  \code{model\_reader} \\
\code{app.authz.writer-role} & \code{ARCHIMESH\_AUTHZ\_WRITER\_ROLE} &
  \code{model\_writer} \\
\code{quarkus.oidc.enabled} & \code{ARCHIMESH\_QUARKUS\_OIDC\_ENABLED} &
  \code{false} \\
\bottomrule
\end{longtable}
\endgroup
